% ============================================================
% General Conclusion
% ============================================================

\chapter*{General Conclusion}
\addcontentsline{toc}{chapter}{General Conclusion}
\label{chap:general-conclusion}

% ---- Summary of achieved work --------------------------------
\section*{Summary of Achieved Work}

This thesis has presented the design, implementation, and pilot validation of an
IoT-based monitoring and deterministic agronomic analytics platform for alfalfa
cultivation in Saharan agricultural environments. The work spanned hardware
architecture, embedded firmware, data communication, local and remote storage,
backend engineering, dashboard development, deterministic analytics, and an
AI-assisted interpretation interface.

The implemented system comprises three embedded sensor nodes communicating over
a 433~MHz LoRa wireless link within a structured ten-minute superframe. The MAIN
gateway node coordinates reception, maintains an authoritative real-time clock, and
persists all acquired measurements on an SD card independent of network availability.
Node2 acquires soil moisture, temperature, and electrical conductivity at a second
irrigation pivot via RS485/Modbus; Node3 acquires atmospheric temperature, humidity,
and barometric pressure. Eight firmware operating modes address the specific field
constraints of the deployment: Power Saving Mode reduces energy consumption through
deep sleep and relay-controlled sensor power; Recovery Mode corrects timing drift
that caused persistent communication failure during prototype testing; Upload Mode
and RTC Synchronisation Mode manage the interface between offline field operation
and online data transfer.

The backend, built with Node.js and PostgreSQL, enforces a strict timestamp policy
that separates the RTC-derived measurement timestamp (\texttt{measured\_at}) from
the upload receipt time, which is stored only as transfer audit metadata. A
React-based dashboard presents soil and weather trends, agronomic events, and the
output of a ten-stage deterministic quality-control and reliability-scoring pipeline.
The Processed Data Viewer exposes the full analytical lineage in read-only form.
AI-assisted interpretation is provided through a backend proxy that enforces six
boundary rules preventing Gemini from computing primary analytics or modifying
stored data.

% ---- Response to the problem statement -----------------------
\section*{Response to the Problem Statement}

Each of the five problems identified in the general introduction has been addressed
by the implemented system.

The absence of a local storage mechanism for offline operation is addressed by
the SD-based Local Storage Mode, which writes sensor readings and system events
to the SD card at the end of every measurement cycle, independently of network
availability. Upload Mode transfers the accumulated records in a single batch
session triggered by the operator.

The low-power constraint is addressed by Power Saving Mode, which combines deep
sleep between measurement cycles with relay-controlled sensor power and BME280
forced measurement mode on Node3, eliminating continuous current draw from both
the microcontrollers and the sensors.

The timestamp integrity problem is addressed at every layer of the system. The
DS3231 RTC provides the canonical \texttt{measured\_at} timestamp, written by
MAIN at packet reception and never overwritten. The backend enforces this
separation at ingestion, and all dashboard charts draw their time axes from
\texttt{measured\_at}.

The absence of a deterministic analytics layer is addressed by the ten-stage
QC and reliability-scoring pipeline, which produces auditable analytics snapshots
that are authoritative for all dashboard metrics. Gemini receives only these
snapshots and is explicitly excluded from all computation stages.

The alfalfa monitoring gap is addressed by deploying the prototype specifically
for alfalfa cultivation in the El~Oued region, collecting the soil and atmospheric
measurements relevant to this crop's water, temperature, and salinity requirements
in a Saharan context.

% ---- Technical contributions ---------------------------------
\section*{Technical Contributions}

The primary contributions of this work are: an offline-first three-node LoRa
architecture with deterministic superframe coordination; a timestamp traceability
policy grounded in RTC hardware; eight documented firmware operating modes each
motivated by a specific field constraint; a ten-stage deterministic analytics
pipeline with an explicit AI boundary; a Processed Data Viewer for analytical
transparency; and a documented pilot deployment targeting alfalfa monitoring in
Saharan agriculture.

% ---- Validation summary (cautious) ---------------------------
\section*{Validation Summary}

The system was validated over a pilot window from 2026-05-06 to 2026-05-11.
The validation evidence comprises 2433 sensor readings distributed equally across
three nodes, eight system events, six manually recorded agronomic events,
six upload sessions all reaching processed status, and 135 rows of deterministic
pipeline output. These results confirm the functional feasibility of the
end-to-end architecture within the scope of a pilot deployment. They do not
constitute full-season agronomic validation, and the reliability indicator of
0.399 is presented as a pilot-phase output of the uncalibrated pipeline.

% ---- Limitations ---------------------------------------------
\section*{Limitations}

The system operates within a well-defined scope. Soil electrical conductivity
is treated as a relative trend indicator; raw in-situ EC values are not converted
to official salinity classifications. The sensor stack does not include pH or NPK
measurement. Evapotranspiration cannot be computed because the weather node lacks
wind speed and solar radiation sensors. The validation window covers five-and-a-half
days, not a complete alfalfa growth cycle. Battery autonomy and LoRa communication
range are design intentions that were not quantitatively characterised during the
pilot deployment. The Pivot~1 local soil measurement path requires register map
confirmation before validated agronomic data can be collected from that sensor.
The effects of precipitation and high humidity on LoRa communication reliability
and enclosure performance were not experimentally isolated during the pilot
deployment and represent an open environmental validation question.

% ---- Future work ---------------------------------------------
\section*{Future Work}

Several directions for future development emerge from the present work.

\textbf{Extended validation.} The most immediate priority is a deployment covering
one or more complete alfalfa cut cycles to assess seasonal measurement coverage,
long-term node reliability, and the agronomic interpretability of multi-week
soil moisture trends.

\textbf{Battery autonomy study.} Instrumenting the nodes with battery voltage
sensing hardware would allow Power Saving Mode to be characterised in quantitative
terms under actual field discharge conditions.

\textbf{LoRa range characterisation.} A structured range trial with RSSI and SNR
logging at known node-to-gateway distances would establish the communication
envelope of the system and support placement decisions for future deployments.

\textbf{Sensor calibration.} Local calibration of the RS485 soil sensors against
gravimetric soil moisture samples and EC laboratory measurements would allow the
system to report calibrated agronomic values rather than relative trend indicators.

\textbf{Reliability threshold calibration.} The reliability scoring pipeline
requires field calibration against known-good and known-poor measurement windows.
Once calibrated, the 0.399 pilot-phase indicator can be replaced by a threshold-
grounded score that supports operator alerts with a defined confidence level.

\textbf{Agronomic expert validation.} Collaboration with an agronomist specialising
in alfalfa cultivation in arid environments would establish whether the soil and
atmospheric measurements collected by the system are sufficient to support
irrigation scheduling recommendations, and what additional data points would be
required.

\textbf{Mobile upload workflow.} Replacing the push-button upload trigger with a
mobile application would reduce the operator effort required for periodic field
visits and enable upload sessions without a laptop at the field site.

\textbf{Crop and region expansion.} The monitoring platform is not inherently
specific to alfalfa or El~Oued. Deploying the same architecture for other
drought-sensitive crops grown in Saharan conditions — such as date palms, onions,
or tomatoes under drip irrigation — would validate the generalisability of the
design and provide a basis for comparative agronomic analysis across crop types.
These extensions should be preceded by appropriate agronomic consultation and
sensor calibration specific to each crop and site.
